\section{Materials and Methods}\label{sec:methods}

This section specifies the experimental data, the dual-head classifier
shared by every architecture in the study, the proposed Custom EfficientNet
V2 backbone, the two AttentionHub variants, the masked multi-task loss, and
the training, evaluation, and ablation protocols. Throughout, the design is
governed by a single principle: every architectural choice that is not
inherited verbatim from a published reference is later subjected to a
controlled measurement, so that the proposed model is the conclusion of an
ablation rather than an unverified assertion.

\subsection{Dataset}\label{ssec:dataset}

The experimental data combine two complementary, publicly released oral
lesion image collections that together support the dual-task formulation.
\emph{Dataset~1} (DS1) is the \emph{Oral Images Dataset} of
Chandrashekar~\textit{et~al.}~\cite{chandrashekar2021oral} (Mendeley Data,
\textsc{cc by}~4.0), a curated set of 165 benign and 158 malignant oral
lesion photographs collected in consultation with clinicians at hospitals
and colleges across Karnataka, India; it drives the binary head of the
classifier. \emph{Dataset~2} (DS2) is the \emph{Mouth and Oral Diseases}
(MOD) dataset of Rashid~\textit{et~al.}~\cite{rashid2022mod,ref19}
(distributed via Kaggle), compiled from dental clinics in Okara, Punjab,
Pakistan and supplementary online sources, providing seven-way subtype
supervision across the classes \textbf{CaS}, \textbf{CoS}, \textbf{Gum},
\textbf{MC}, \textbf{OC}, \textbf{OLP}, and \textbf{OT}, with labels
contributed by expert dental practitioners; we retain the source
dataset's class abbreviations verbatim rather than assigning clinical
expansions that the release does not document. Within DS2, the subset
$\{\text{MC},\text{OC},\text{CaS}\}$ is taken to denote the malignant
subtypes, and the remaining four classes the non-malignant subtypes, for
the purpose of constructing binary supervision when DS2 images are read
into the binary head. Representative images for the two binary categories
and the seven subtype classes are shown in Fig.~\ref{fig:dataset}; the
visual overlap in colour, texture, and boundary cues across several
subtypes motivates the multi-class formulation pursued here.

The original DS2 release exhibited a distributional similarity between its
provided \emph{Training} and \emph{Testing} folders that would inflate
test-set scores by leaking near-duplicate appearance statistics across the
split. To eliminate this confound, the DS2 \emph{Training} and
\emph{Validation} directories are merged and re-partitioned with a
stratified 60/20/20 train/validation/test split, seeded for reproducibility
(seed${}={}$42). DS1 is partitioned independently with a stratified
80/10/10 split. The resulting test set contains 1\,646 DS2 images
distributed across the seven subtype classes, with the following per-class
supports: CaS${}={}$256, CoS${}={}$239, Gum${}={}$192, MC${}={}$288,
OC${}={}$173, OLP${}={}$288, and OT${}={}$210. The full per-class
distribution across the merged train, validation, and test partitions is
plotted in Fig.~\ref{fig:classdist}.

All images are resized to $224 \times 224$ and normalised with ImageNet
channel statistics. Training augmentation comprises random horizontal and
vertical flips, random rotation up to $\pm 15^{\circ}$, a random crop taken
after resizing to $256 \times 256$, and colour jitter (brightness,
contrast, and saturation each $\pm 0.2$; hue $\pm 0.1$). The validation and
test sets receive no augmentation. The test set is held out in its entirety
and consulted only for final reporting, so that no design decision in this
paper is informed by test-set performance.

\begin{figure*}[tbp]
  \centering
  \includegraphics[width=\textwidth,height=0.92\textheight,%
  keepaspectratio]{fig02_dataset_sample_grid}
  \caption{Representative images from the two datasets, arranged as a
  9-row by 4-column grid: one class per row, with four randomly sampled
  examples per class shown across the four columns (seed${}={}$42). The top
  two rows show the DS1 binary categories (benign, malignant); the lower
  seven rows show the DS2 subtype classes CaS, CoS, Gum, MC, OC, OLP, and
  OT. Several subtypes share overlapping colour, texture, and boundary
  appearance, which is the visual difficulty that the seven-class head must
  resolve.}
  \label{fig:dataset}
\end{figure*}

\begin{figure}[!t]
  \centering
  \includegraphics[width=\columnwidth]{fig02b_class_distribution}
  \caption{Class distribution across the two datasets. \emph{Left:} DS1
  benign-versus-malignant image counts. \emph{Right:} for each DS2 subtype,
  the total number of images across the merged train, validation, and test
  partitions, alongside the held-out test-set support used for per-class
  reporting in Section~\ref{sec:results}.}
  \label{fig:classdist}
\end{figure}

\subsection{Dual-Head Multi-Task Classifier}\label{ssec:dualhead}

All ten architectures evaluated in this study (nine baseline backbones
and the proposed model) share a common dual-head wrapper, denoted
\texttt{MultiTaskOralClassifier}. A shared backbone produces a pooled
feature vector, after which two independent multilayer perceptron (MLP)
heads predict the binary class (2 logits) and the subtype class (7 logits).
Each head is a \texttt{Linear}$(F\!\rightarrow\!512)\rightarrow
\text{ReLU}\rightarrow\text{Dropout}\rightarrow
\texttt{Linear}(512\!\rightarrow\!C)$ stack, where $F$ is the
backbone feature dimensionality and $C$ is the head's class count. A shared
dropout layer with rate $p=0.5$ is applied to the pooled features before
either head, providing the principal regularisation against over-fitting on
a dataset of modest scale.

The dual-head construction is what allows DS1, which carries only binary
labels, and DS2, which carries both binary and subtype labels, to train a
single model jointly: the subtype loss is masked for samples that lack a
subtype label, as detailed in Section~\ref{ssec:loss}. Sharing the backbone
across both tasks also forces the learned representation to be useful for
benign-versus-malignant screening and for fine-grained subtype
discrimination simultaneously, rather than specialising to either task in
isolation.

\subsection{Custom EfficientNet V2 Backbone}\label{ssec:backbone}

The proposed backbone is a five-stage CNN derived from the
\texttt{tf\_efficientnetv2\_b0} reference, the smallest member of the
EfficientNetV2 family~\cite{tan2021effv2}. Two structural
changes are applied to the donor architecture, and only these two; every
other stage is retained verbatim from the EfficientNetV2-B0 reference so
that the contribution of each change can be measured cleanly.

\begin{enumerate}
  \item \textbf{Stage-4 replacement.} The donor's Block-4, a Mobile
  Inverted Bottleneck Convolution (MBConv) block fused with a
  Squeeze-and-Excitation (SE) unit mapping $96 \rightarrow 112$ channels,
  is removed and replaced by an \emph{AttentionHub}
  (Section~\ref{ssec:hub}) with identical input and output channel counts.
  The hub is the only stage whose internal structure departs from the
  EfficientNetV2-B0 reference, and it is the only component the ablation of
  Section~\ref{sec:ablation} varies.
  \item \textbf{Tail trimming.} The donor's Block-6 and its
  \texttt{conv\_head} are dropped as redundant for this task. The
  192-channel output of Stage~5 is pooled directly into the dual-head
  classifier, which removes a substantial parameter and compute cost
  carried by the donor's wide convolutional head.
\end{enumerate}

The resulting backbone comprises five stages: \emph{Stem${}+{}$Block-0
${}+{}$Block-1} (32 channels) $\rightarrow$ \emph{Block-2} (48)
$\rightarrow$ \emph{Block-3} (96) $\rightarrow$ \emph{AttentionHub} (112)
$\rightarrow$ \emph{Block-5} (192), followed by global average pooling
and the two heads. The complete model contains between 4.79~M and 4.80~M
parameters and requires between 0.493 and 0.495 GFLOPs across the
attention-bearing ablation variants reported in Section~\ref{sec:ablation};
the no-attention control, which restores the donor's heavier
MBConv${}+{}$SE Block-4, is the exception at 5.70~M parameters and 0.636
GFLOPs (Table~\ref{tab:ablation}). The proposed configuration sits at the
lower end of the attention-bearing range (4.79~M parameters, 0.493 GFLOPs).
The full backbone, the Stage-4 hub, the global average pooling stage, and
the dual-head split are diagrammed in Fig.~\ref{fig:arch}.

\begin{figure*}[!t]
  \centering
  \includegraphics[width=\textwidth]{fig01_custom_efficientnet_v2_arch}
  \caption{The proposed Custom EfficientNet V2. A five-stage backbone
  derived from EfficientNetV2-B0 feeds a global average pooling layer and a
  shared dropout, after which two independent MLP heads predict the binary
  (benign / malignant) and subtype (7-class) labels. Stage~4 is replaced by
  the proposed AttentionHub-v2, a sequential Triplet$\rightarrow$SE cascade
  flanked by $1\times1$ channel-reduction and channel-expansion
  projections; the donor's Block-6 and convolutional head are removed. The
  complete model has 4.79~M parameters and 0.493 GFLOPs.}
  \label{fig:arch}
\end{figure*}

\subsection{AttentionHub Variants}\label{ssec:hub}

The AttentionHub is the single module whose composition changes across
ablation cells and across the two versions of the design presented here.
Both versions occupy Stage~4 of the backbone and preserve the
$96 \rightarrow 112$ channel mapping of the block they replace.

\paragraph{AttentionHub-v1 (parallel triple)}
The originally proposed hub routes the Stage-3 output through three parallel
attention branches: Bottleneck Attention Module (BAM), Triplet Attention,
and Kolmogorov--Arnold Network attention (KAN-Attention). Each branch is preceded
by a $1\times1$ channel-reduction projection that halves the branch channel
count. The three branch outputs are concatenated and fused by a $1\times1$
projection back to 112 channels. Ablation of v1 is implemented by removing
branches: the empty configuration restores the donor's original Block-4
(the canonical no-attention control), and any subset of
$\{\text{BAM},\text{Triplet},\text{KAN}\}$ is admissible. The three
candidate modules, whose internal structure is shown in
Fig.~\ref{fig:modules}, are as follows.

\begin{itemize}
  \item \textbf{BAM}~\cite{park2018bam} combines a channel gate (global
  average pooling followed by an MLP bottleneck) and a spatial gate
  ($1\times1 \rightarrow$ dilated $3\times3 \rightarrow 1\times1$, producing
  a single-channel spatial map). The two gates are summed, passed through a
  sigmoid, and multiplied with the input. BAM therefore exercises
  \emph{both} a channel and a spatial role.
  \item \textbf{Triplet Attention}~\cite{misra2021triplet} uses three
  branches, each applying a $7\times7$ attention gate to a permutation of
  the input tensor: the first branch swaps the channel and height axes, the
  second swaps the channel and width axes, and the third leaves the input
  in its native $(B,C,H,W)$ layout. Each gate applies a Z-Pool, which
  concatenates max- and mean-pooled statistics across the gated axis into
  two channels, followed by a $7\times7$ convolution and a sigmoid. The
  three branch outputs are averaged. Triplet exercises a
  \emph{cross-dimensional spatial} role at near-zero additional parameter
  cost.
  \item \textbf{KAN-Attention} is a Kolmogorov--Arnold-Network-inspired
  channel recalibration~\cite{liu2024kan}. Global average pooling produces
  a per-channel scalar; each scalar is passed through a learnable B-spline
  basis expansion (five Gaussian basis functions on $[-3,+3]$) with a SiLU
  residual, then sigmoid-gated and multiplied with the input. KAN-Attention
  is \emph{purely channel-wise} and supplies a non-linear alternative to the
  linear MLP used in the SE and BAM channel gates.
\end{itemize}

\paragraph{AttentionHub-v2 (sequential Triplet$\rightarrow$SE)}
The proposed model uses a sequential cascade in which the channel-reduced
feature is first refined by Triplet Attention and then by
Squeeze-and-Excitation (SE)~\cite{hu2018se}. The SE module pools globally,
passes the pooled descriptor through a two-layer $1\times1$ bottleneck with
reduction ratio 16 and a sigmoid gate, and multiplies the result back into
the input. No LayerScale, no per-module residual, and no additional gating
are introduced: both attentions are multiplicative gates that preserve the
input information without an explicit skip connection. The composition
order, spatial attention before channel attention, follows the
convention established by the Convolutional Block Attention Module
(CBAM)~\cite{woo2018cbam}. The internal layout of the v2 hub is shown in
Fig.~\ref{fig:hubv2}, and a side-by-side contrast with the v1 parallel hub
is given in Fig.~\ref{fig:v1v2}. The design rationale for v2 (in
particular, why a sequential cascade is preferred over the v1 parallel
fusion and why SE is chosen as the channel partner) is not asserted here
but \emph{derived} in Section~\ref{sec:ablation} from the v1 ablation.

\begin{figure*}[!t]
  \centering
  \includegraphics[width=\textwidth]{fig01d_attention_modules_internals}
  \caption{Internal structure of the four attention modules used in this
  study. \emph{Triplet Attention} applies a $7\times7$ gate to three
  axis-permuted views of the feature map and averages them, exercising a
  cross-dimensional spatial role. \emph{Squeeze-and-Excitation (SE)}
  recalibrates channels through a global-average-pooling bottleneck.
  \emph{BAM} sums a channel gate and a spatial gate, exercising a mixed
  role. \emph{KAN-Attention} replaces the linear channel MLP with a
  learnable B-spline basis expansion, remaining purely channel-wise.}
  \label{fig:modules}
\end{figure*}

\begin{figure}[!t]
  \centering
  \includegraphics[width=\columnwidth]{fig01b_attentionhub_v2_detail}
  \caption{The proposed AttentionHub-v2. A $1\times1$ projection reduces the
  96-channel Stage-3 feature; Triplet Attention then applies
  cross-dimensional spatial gating, SE applies pure per-channel
  recalibration, and a $1\times1$ projection expands the result to 112
  channels for Stage~5. The two attention modules occupy disjoint functional
  roles (spatial and channel), which Section~\ref{sec:ablation} shows is
  the condition required to avoid the observed sub-baseline accuracy
  ceiling.}
  \label{fig:hubv2}
\end{figure}

\begin{figure*}[!t]
  \centering
  \includegraphics[width=\textwidth]{fig01c_attentionhub_v1_vs_v2}
  \caption{AttentionHub-v1 versus AttentionHub-v2. \emph{Left:} the v1 hub
  fuses BAM, Triplet, and KAN-Attention in parallel, concatenating their
  channel-reduced branch outputs before a $1\times1$ projection.
  \emph{Right:} the proposed v2 hub composes Triplet Attention and SE
  sequentially. The two hubs reuse the same building blocks but differ in
  composition; Section~\ref{sec:ablation} shows that the difference in
  composition is what separates their accuracy.}
  \label{fig:v1v2}
\end{figure*}

\subsection{Multi-Task Loss}\label{ssec:loss}

The two heads are trained jointly under a single combined cross-entropy
(CE) objective,

\begin{equation}\label{eq:loss}
  \begin{split}
  \mathcal{L} \;=\;{}& w_b \cdot \mathrm{CE}\!\left(p_{\text{binary}},\,
  y_{\text{binary}}\right) \\
  &{}+\; w_s \cdot \mathrm{CE}\!\left(
  p_{\text{subtype}},\, y_{\text{subtype}};\,
  \texttt{ignore\_index}=-1\right),
  \end{split}
\end{equation}

with equal task weights $w_b = w_s = 1$. The \texttt{ignore\_index}${}={}-1$
convention is the mechanism that lets the two datasets co-train against one
model: DS1 samples, which carry no subtype label, are assigned the sentinel
subtype label $-1$ and therefore contribute only to the binary term, while
DS2 samples contribute to both terms. A numerical guard converts any
not-a-number (NaN) subtype loss to zero, so that a mini-batch composed
exclusively of DS1 samples, for which the subtype term is undefined,
remains a valid training step.

The joint update is summarised in Algorithm~1, typeset in
Fig.~\ref{alg:train}. The masked subtype loss is the only coupling between
the binary-only DS1 and the dually labelled DS2: DS1 samples shape the
shared backbone solely through $\mathcal{L}_b$, whereas DS2 samples shape it
through both $\mathcal{L}_b$ and $\mathcal{L}_s$.

\begin{figure}[!t]
  \centering
  \fbox{%
  \begin{minipage}{0.92\columnwidth}
  \ttfamily \footnotesize
  \textbf{Algorithm 1:} Joint masked multi-task training step.\\[2pt]
  \textbf{Input:} mini-batch $(x, y_b, y_s)$; model $M = $ backbone\\
  \hphantom{\textbf{Input:} }${}+ (g_{\text{binary}}, g_{\text{subtype}})$;
  optimizer Opt;\\
  \hphantom{\textbf{Input:} }loss weights $(w_b, w_s)$.\\
  \textbf{Output:} updated model parameters.\\[3pt]
  1.\ \ $F \leftarrow$ backbone$(x)$ \hfill \# $B \times D$\\
  2.\ \ $F \leftarrow$ Dropout$(F,\ p = 0.5)$\\
  3.\ \ $\hat{y}_b \leftarrow g_{\text{binary}}(F)$ \hfill \# $B \times 2$\\
  4.\ \ $\hat{y}_s \leftarrow g_{\text{subtype}}(F)$ \hfill \# $B \times 7$\\
  5.\ \ $L_b \leftarrow$ CE$(\hat{y}_b,\ y_b)$ \hfill \# always defined\\
  6.\ \ $L_s \leftarrow$ CE$(\hat{y}_s,\ y_s;$ ignore\_index $= -1)$\\
  \hphantom{6.\ \ }\hfill \# NaN if every $y_s = -1$\\
  7.\ \ \textbf{if} isnan$(L_s)$\textbf{:}\ \ $L_s \leftarrow 0$
  \hfill \# DS1-only batch guard\\
  8.\ \ $L_{\text{total}} \leftarrow w_b \cdot L_b + w_s \cdot L_s$\\
  9.\ \ Opt.zero\_grad()\textbf{;}\ \ $L_{\text{total}}$.backward()%
  \textbf{;}\ \ Opt.step()
  \end{minipage}}
  \caption{Joint training step under the masked multi-task loss. The
  sentinel subtype label $-1$ marks DS1 (binary-only) samples; the
  \texttt{ignore\_index} convention and the NaN guard on line~7 together
  allow DS1 and DS2 to co-train a single shared backbone.}
  \label{alg:train}
\end{figure}

\subsection{Training Protocol}\label{ssec:training}

All ten models (the nine baseline backbones and the proposed model)
are trained from scratch under a single matched recipe, deliberately chosen
to guarantee a fair comparison. \textbf{No ImageNet pre-trained weights are
used.} Pre-training would smuggle a representational prior into the
comparison and confound the contribution of the dataset itself with the
contribution of the architecture; training every model from random
initialisation removes that confound and isolates what each backbone learns
from the oral-disease data alone. The recipe is fixed as follows.

\begin{itemize}
  \item \textbf{Optimiser:} Adam~\cite{kingma2015adam}, learning rate
  $1\times10^{-4}$, weight decay $1\times10^{-4}$.
  \item \textbf{Scheduler:} cosine annealing over the full epoch budget.
  \item \textbf{Batch size:} 64, with gradient accumulation factor 1
  (effective batch size 64).
  \item \textbf{Epoch budget:} 200, with early stopping (patience 15,
  minimum improvement $\Delta = 1\times10^{-4}$ on validation loss).
  \item \textbf{Seed:} 42, with deterministic CuDNN.
  \item \textbf{Image size:} $224\times224$.
  \item \textbf{Hardware:} a single NVIDIA GeForce RTX~4060~Ti.
\end{itemize}

This recipe is enforced for every model without exception. In particular,
the AttentionHub ablation variants are forced onto the identical baseline
recipe, so that v1, v2, and the seven ablation cells of
Section~\ref{sec:ablation} are directly comparable. Per-model
hyper-parameter tuning is deliberately \emph{not} performed: tuning would
advantage the models we elect to tune and disadvantage those we do not,
reintroducing exactly the kind of uncontrolled variable that a fair
benchmark must eliminate. Where a baseline encounters memory pressure on the
16~GB GPU, the correct response under this protocol is gradient
accumulation at a fixed effective batch size, not a per-model
hyper-parameter override, a stance established by prior fair-comparison
work and retained here as a methodological commitment rather than a
convenience.

\subsection{Evaluation Protocol}\label{ssec:evaluation}

For every model we report the following quantities on the held-out test
set.

\begin{itemize}
  \item \textbf{Classification metrics:} accuracy, weighted precision,
  weighted recall, and weighted F1, for both the binary head (the DS1 test
  split together with the full DS2 test set, $1\,711$ images in total) and the
  subtype head (the DS2 seven-class test set, $1\,646$ images).
  \item \textbf{Per-class report:} precision, recall, F1, and support for
  each of the seven subtype classes.
  \item \textbf{Confusion matrices:} the binary and subtype confusion
  matrices, reported side by side.
  \item \textbf{Computational metrics:} parameter count, GFLOPs computed
  with \texttt{thop}, on-disk model size in megabytes, GPU peak memory
  during inference,\footnote{Compute is reported as the multiply-accumulate
  (MAC) count returned by the \texttt{thop} profiler, following common
  practice; we use the label GFLOPs for these MAC-count values for
  consistency with the comparison literature.} the mean and standard
  deviation of single-image
  inference latency (a 200-iteration warm-up followed by 1\,000 timed
  iterations), the P50 / P90 / P95 / P99 latency percentiles (of which
  the P50 and P95 values are tabulated in Section~\ref{sec:results}),
  energy in kWh and carbon in kg~CO\textsubscript{2}-equivalent measured
  with \texttt{codecarbon} where available, total training time, and the
  number of epochs completed before early stopping.
  \item \textbf{Explainability artifacts:} for each model, a
  Gradient-weighted Class Activation Mapping++ (Grad-CAM++) panel and a Local
  Interpretable Model-Agnostic Explanations (LIME) boundary-mask panel for
  the binary head, and the matching pair for the subtype head, evaluated on
  a class-balanced sample of test images.
\end{itemize}

The per-model JSON metric files, training-time logs, latency histograms,
confusion matrices, and explainability images are released alongside the
paper so that every reported number can be independently regenerated.

\subsection{Ablation Protocol}\label{ssec:ablation_protocol}

The AttentionHub-v1 ablation is implemented as a code-level branch switch:
each ablation key maps to a tuple of active branches, a subset of
$\{\text{bam},\text{triplet},\text{kan}\}$, where the empty tuple falls
back to the donor's original Block-4 and thus serves as the no-attention
control. Because the switch operates inside the hub and changes nothing
else in the network or the recipe, the eight cells differ from one another
only in which attention branches are active, the property required for
the ablation to be causally interpretable. The eight cells are listed in
Table~\ref{tab:ablkeys}, together with the functional role each
configuration's hub exercises.

\begin{table}[!t]
  \centering
  \caption{AttentionHub-v1 ablation keys. Each key activates a subset of the
  three candidate branches; the empty configuration restores the donor
  MBConv${}+{}$SE Block-4 as the no-attention control. The role column
  classifies the configuration by the functional role its hub exercises and
  is the basis of the role-complementarity analysis in
  Section~\ref{sec:ablation}.}
  \label{tab:ablkeys}
  \begin{tabularx}{\columnwidth}{@{}l l X@{}}
    \hline
    Key & Active branches & Functional role \\
    \hline
    \texttt{none} & --- (donor MBConv${}+{}$SE) & no-attention control \\
    \texttt{bam} & BAM & spatial${}+{}$channel singleton \\
    \texttt{triplet} & Triplet & cross-dim-spatial singleton \\
    \texttt{kan} & KAN & channel singleton \\
    \texttt{bam\_triplet} & BAM${}+{}$Triplet & two-spatial-role pair \\
    \texttt{bam\_kan} & BAM${}+{}$KAN & mixed-role pair \\
    \texttt{triplet\_kan} & Triplet${}+{}$KAN &
      cross-dim-spatial${}+{}$pure-channel pair \\
    \texttt{full} (${}={}$v1) & BAM${}+{}$Triplet${}+{}$KAN &
      original proposed model \\
    \hline
  \end{tabularx}
\end{table}

In addition, we report a v2-EMA run, a sequential Triplet${}+{}$Efficient
Multi-scale Attention (EMA) cascade, which is \emph{not} an ablation cell
but a deliberate stress test of the role-complementarity principle: it
replaces SE with EMA, a module that combines a multi-scale spatial branch
with a channel branch. Because EMA exercises a spatial role, the principle
predicts that the v2-EMA cascade should regress, and the result reported in
Section~\ref{sec:ablation} confirms this prediction. Every ablation cell is
trained under the identical baseline recipe used by the nine baseline
backbones; no recipe exceptions apply.
